% !TeX root = ../../../main.tex
\ChFiveAbstract
{
Clinical narratives encode temporal dynamics essential for modeling patient trajectories, yet large-scale temporally annotated resources are scarce. We introduce PMOA-TTS, a corpus of 124{,}349 single-patient PubMed Open Access case reports converted into structured textual timelines of (event, time) pairs using a scalable large-language-model pipeline (\texttt{Llama 3.3 70B} and \texttt{DeepSeek-R1}). The corpus comprises over 5.6 million timestamped events, alongside extracted demographics and diagnoses. Technical validation uses a clinician-curated gold set and three measures: semantic event matching, temporal concordance (c-index), and alignment error summarized with Area Under the Log-Time CDF (AULTC). We benchmark alternative prompting and model choices and provide documentation to support reproduction. 
\ChFiveRebuttal{
PMOA-TTS provides a large, publicly available resource for developing and evaluating methods for clinical timeline extraction and temporal reasoning from narrative text.
}
% PMOA–TTS enables research on timeline extraction, temporal reasoning, survival modeling and event forecasting from narrative text, and offers broad diagnostic and demographic coverage. 
Data and code are openly available in public repositories.
% PMOA–TTS dataset is hosted on Hugging Face Datasets at \url{https://huggingface.co/datasets/snoroozi/pmoa-tts}. Our full codebase, including preprocessing scripts, prompt templates, and evaluation pipelines,
% is available at: \url{https://github.com/jcweiss2/pmoa_tts}. 
}.

%\ChFiveKeywords{keyword1, Keyword2, Keyword3, Keyword4}
\ChFiveKeywords{clinical natural language processing, temporal information extraction, textual time series, large language models, PubMed Open Access case reports}



% Background & Summary
% Methods
% Data Records
% Data Overview (optional)
% Technical Validation
% Usage Notes (optional)
% Data Availability
% Code Availability
% References
% Author Contributions
% Competing Interests
% Acknowledgements (optional)
% Funding
% Ethics statement* 

%%%%%%%%%%%%%%%%%%%%%%%%%%%%%%%%%%%%%%%%%%%%%%
%%%%%%%%%%%%%% Introduction %%%%%%%%%%%%%%%%%
%%%%%%%%%%%%%%%%%%%%%%%%%%%%%%%%%%%%%%%%%%%%%
